% GENERATED FILE. Edit canonical lessons, then run npm run generate:books.
% canonical-source-hash: fnv1a64:26585fdbba359acf
% canonical-lessons: FR-C522-poesie, FR-C522-bande-dessinee, FR-C522-polar, FR-C522-biographie, FR-C522-conte

\chapter{Quel genre de livres ?}
\label{ch:fr-quel-genre-de-livres}
\hlchaptermodality{522}

\begin{chapteropening}
\textbf{By the end of this chapter:} \emph{I can say what kinds of books I like.}

Complete the last lesson of chapter 522: I can say what kinds of books I like.
\end{chapteropening}

\section[la poésie]{la poésie — poetry}

\label{lesson:FR-C522-poesie}



\begin{quote}
Before the new one: say the French for signed message, then the French for to sign.
\end{quote}

\subsection*{You'll want to know: la poésie}
\textbf{la poésie} — \textquotedblleft{}poetry\textquotedblright{}.

\begin{quote}
Mon père lit de la poésie le soir.
\end{quote}

\begin{quote}
\emph{My father reads poetry in the evening.}
\end{quote}

\emph{Un poème} is a single poem.

\begin{cousinweb}
\emph{la poésie} comes from Latin \emph{poēsis} (EtymDB), a borrowing from Greek \emph{poíēsis} \textquotedblleft{}a making, a poem\textquotedblright{}.
\end{cousinweb}

\begin{grammarlens}[title={What you've built}]
One more word for telling, explaining and describing at B1.
\end{grammarlens}

\subsection*{Your turn}
\begin{itemize}
\raggedright
  \item \emph{Say it:} \emph{la poésie}
  \item \emph{Say it:} \emph{la poésie}, once more
  \item \emph{Say it:} a sentence of your own with \emph{la poésie}
  \item \emph{Recall:} say the French for signed message, then the French for to sign, then say \emph{la poésie} again
\end{itemize}

\begin{tcolorbox}[breakable,colback=teal!4,colframe=teal!35!black,title={Before you move on}]
What does la poésie mean? (\textquotedblleft{}Poetry\textquotedblright{}.) Say it once more.
\end{tcolorbox}
\section[la bande dessinée]{la bande dessinée — comic book, comic strip}

\label{lesson:FR-C522-bande-dessinee}



\begin{quote}
Before the new one: say the French for to sign, then the French for poetry.
\end{quote}

\subsection*{You'll want to know: la bande dessinée}
\textbf{la bande dessinée} — \textquotedblleft{}comic book, comic strip\textquotedblright{}.

\begin{quote}
Les adultes aussi lisent des bandes dessinées.
\end{quote}

\begin{quote}
\emph{Adults read comics too.}
\end{quote}

Informally \emph{la BD}. A big art form in France and Belgium.

\begin{cousinweb}
\emph{la bande dessinée} is \textquotedblleft{}a drawn strip\textquotedblright{}: \emph{bande} is of Germanic origin, and \emph{dessiner} comes from Italian \emph{disegnare}, from Latin \emph{dēsignāre} \textquotedblleft{}to mark out\textquotedblright{}, built on \emph{signum} \textquotedblleft{}a sign\textquotedblright{} (Lewis \& Short).
\end{cousinweb}

\begin{grammarlens}[title={What you've built}]
One more word for telling, explaining and describing at B1.
\end{grammarlens}

\subsection*{Your turn}
\begin{itemize}
\raggedright
  \item \emph{Say it:} \emph{la bande dessinée}
  \item \emph{Say it:} \emph{la bande dessinée}, once more
  \item \emph{Say it:} a sentence of your own with \emph{la bande dessinée}
  \item \emph{Recall:} say the French for to sign, then the French for poetry, then say \emph{la bande dessinée} again
\end{itemize}

\begin{tcolorbox}[breakable,colback=teal!4,colframe=teal!35!black,title={Before you move on}]
What does la bande dessinée mean? (\textquotedblleft{}Comic book, comic strip\textquotedblright{}.) Say it once more.
\end{tcolorbox}
\section[le polar]{le polar — crime novel, detective story (informal)}

\label{lesson:FR-C522-polar}



\begin{quote}
Before the new one: say the French for poetry, then the French for comic book.
\end{quote}

\subsection*{You'll want to know: le polar}
\textbf{le polar} — \textquotedblleft{}crime novel, detective story (informal)\textquotedblright{}.

\begin{quote}
En été, je lis surtout des polars.
\end{quote}

\begin{quote}
\emph{In summer I mostly read crime novels.}
\end{quote}

Informal for \emph{le roman policier}. Also a crime film.

\begin{cousinweb}
\emph{le polar} is a slang shortening of \emph{(roman) policier} \textquotedblleft{}detective (novel)\textquotedblright{} with the familiar ending \emph{-ar} (EtymDB).
\end{cousinweb}

\begin{grammarlens}[title={What you've built}]
One more word for telling, explaining and describing at B1.
\end{grammarlens}

\subsection*{Your turn}
\begin{itemize}
\raggedright
  \item \emph{Say it:} \emph{le polar}
  \item \emph{Say it:} \emph{le polar}, once more
  \item \emph{Say it:} a sentence of your own with \emph{le polar}
  \item \emph{Recall:} say the French for poetry, then the French for comic book, then say \emph{le polar} again
\end{itemize}

\begin{tcolorbox}[breakable,colback=teal!4,colframe=teal!35!black,title={Before you move on}]
What does le polar mean? (\textquotedblleft{}Crime novel, detective story (informal)\textquotedblright{}.) Say it once more.
\end{tcolorbox}
\section[la biographie]{la biographie — biography}

\label{lesson:FR-C522-biographie}



\begin{quote}
Before the new one: say the French for comic book, then the French for crime novel.
\end{quote}

\subsection*{You'll want to know: la biographie}
\textbf{la biographie} — \textquotedblleft{}biography\textquotedblright{}.

\begin{quote}
Je lis la biographie d'un grand musicien.
\end{quote}

\begin{quote}
\emph{I am reading the biography of a great musician.}
\end{quote}

\emph{Une autobiographie} — an autobiography.

\begin{cousinweb}
\emph{la biographie} is a learned word built from Greek \emph{bíos} \textquotedblleft{}life\textquotedblright{} and \emph{gráphein} \textquotedblleft{}to write\textquotedblright{}, by way of Medieval Greek \emph{biographía}.
\end{cousinweb}

\begin{grammarlens}[title={What you've built}]
One more word for telling, explaining and describing at B1.
\end{grammarlens}

\subsection*{Your turn}
\begin{itemize}
\raggedright
  \item \emph{Say it:} \emph{la biographie}
  \item \emph{Say it:} \emph{la biographie}, once more
  \item \emph{Say it:} a sentence of your own with \emph{la biographie}
  \item \emph{Recall:} say the French for comic book, then the French for crime novel, then say \emph{la biographie} again
\end{itemize}

\begin{tcolorbox}[breakable,colback=teal!4,colframe=teal!35!black,title={Before you move on}]
What does la biographie mean? (\textquotedblleft{}Biography\textquotedblright{}.) Say it once more.
\end{tcolorbox}
\section[le conte]{le conte — tale, story}

\label{lesson:FR-C522-conte}



\begin{quote}
Before the new one: say the French for crime novel, then the French for biography.
\end{quote}

\subsection*{You'll want to know: le conte}
\textbf{le conte} — \textquotedblleft{}tale, story\textquotedblright{}.

\begin{quote}
Le soir, la mère lit un conte à ses enfants.
\end{quote}

\begin{quote}
\emph{In the evening the mother reads her children a story.}
\end{quote}

\emph{Un conte de fées} — a fairy tale.

\begin{cousinweb}
\emph{le conte} is built on \emph{conter} \textquotedblleft{}to tell\textquotedblright{}, inherited from Latin \emph{computāre} \textquotedblleft{}to count, to reckon\textquotedblright{}, built on \emph{putāre} (Lewis \& Short): to recount.
\end{cousinweb}

\begin{grammarlens}[title={What you've built}]
One more word for telling, explaining and describing at B1.
\end{grammarlens}

\subsection*{Your turn}
\begin{itemize}
\raggedright
  \item \emph{Say it:} \emph{le conte}
  \item \emph{Say it:} \emph{le conte}, once more
  \item \emph{Say it:} a sentence of your own with \emph{le conte}
  \item \emph{Recall:} say the French for crime novel, then the French for biography, then say \emph{le conte} again
\end{itemize}

\begin{tcolorbox}[breakable,colback=teal!4,colframe=teal!35!black,title={Before you move on}]
What does le conte mean? (\textquotedblleft{}Tale, story\textquotedblright{}.) Say it once more.
\end{tcolorbox}
